% =============================================================================
% 05_ehrenwoertliche_erklaerung.tex — Ehrenwörtliche Erklärung
% =============================================================================
% Dieses Kapitel enthält die Bestätigung der selbständigen Erarbeitung.
% WICHTIG: Gemäss Schulvorgabe muss explizit erwähnt werden, dass die
% Bestimmungen zur Selbständigkeit auch für den Einsatz von
% Künstlicher Intelligenz (KI) als Hilfsmittel gelten.
% =============================================================================

\chapter*{Ehrenwörtliche Erklärung}
\addcontentsline{toc}{chapter}{Ehrenwörtliche Erklärung}

Wir erklären hiermit ehrenwörtlich, dass wir die vorliegende
Interdisziplinäre Projektarbeit (IDPA) selbständig und ohne unerlaubte
fremde Hilfe verfasst haben. Alle verwendeten Quellen und Hilfsmittel
sind vollständig angegeben und im Literaturverzeichnis aufgeführt.

Wörtliche oder sinngemässe Übernahmen aus anderen Werken sind als solche
kenntlich gemacht. Die Arbeit wurde weder in gleicher noch in ähnlicher
Form einer anderen Prüfungsbehörde vorgelegt.

% --- Hinweis zu Künstlicher Intelligenz (gemäss Schulvorgabe) ----------------
% WICHTIG: Dieser Absatz ist gemäss den Bestimmungen der Schule obligatorisch.
Uns ist bekannt, dass die Bestimmungen zur ehrenwörtlichen Erklärung und
zur Eigenständigkeit der Arbeit auch für den Einsatz von Hilfsmitteln in
Form von \textbf{Künstlicher Intelligenz (KI)} gelten. Der Einsatz von
KI-gestützten Werkzeugen (z.\,B.\ Large-Language-Models, KI-basierte
Code-Assistenten oder Text\-generierungs\-dienste) ist gemäss den
geltenden Richtlinien der Schule zu deklarieren und entspricht einer
fremden Hilfe, sofern die daraus resultierenden Inhalte ohne
eigenständige Reflexion und kritische Überarbeitung übernommen werden.

Allfällig eingesetzte KI-Hilfsmittel sowie der Umfang ihrer Verwendung
sind nachfolgend vollständig offengelegt:

\begin{itemize}
  % Füge hier alle verwendeten KI-Werkzeuge ein oder trage «keine» ein.
  \item <<Name des KI-Werkzeugs (z.\,B.\ GitHub Copilot, ChatGPT)>>,
        eingesetzt für: <<Beschreibung des Verwendungszwecks und -umfangs>>
  \item <<weiteres KI-Werkzeug oder «keine weiteren KI-Hilfsmittel»>>
\end{itemize}

\vspace{2cm}

% Unterschriftenfelder — Ort, Datum und Unterschriften aller Autor:innen
\noindent
Ort, Datum: \underline{\hspace{6cm}}

\vspace{1.5cm}
\noindent
\begin{tabular}{p{6cm} p{6cm}}
  \underline{\hspace{5.5cm}} & \underline{\hspace{5.5cm}} \\[0.3cm]
  <<Vorname Nachname>>        & <<Vorname Nachname>>        \\[1.5cm]
  \underline{\hspace{5.5cm}} &                             \\[0.3cm]
  <<Vorname Nachname>>        &                             \\
\end{tabular}
